\subsection{Shared low-rank candidates fail; one LLaVA full-hidden interface passes}
\label{sec:subspace-bottleneck}

The sparse and native-channel tests leave a specific alternative unresolved. The clean-minus-masked hidden-state change may be distributed across individual coordinates while remaining confined to a low-dimensional residual-state subspace shared across examples within each model. We test this possibility before interpreting the complete hidden-state intervention. The experiment fits a separate basis for each model at frozen interfaces: layer 1 at visual positions for Gemma, layer 14 at visual positions for Qwen, and layer 30 at the four answer-adjacent positions for LLaVA.

\subsubsection{No example-shared rank from 1 to 128 passes the bidirectional gate}

For each model, 40 images fit the uncentered basis in Equation~\ref{eq:shared-subspace-projection}; 20 different images select among ranks $1,2,4,8,16,32,64,$ and $128$; and a separate 40-image set is reserved for replication. A candidate must restore and remove between 80\% and 120\% of the mean clean-minus-masked sequence-score change on the selection images. No tested rank passes both mean gates in any model, so no rank advances to strict confirmation. Table~\ref{tab:shared-subspace-summary} reports the largest tested rank on the replication images together with the full-hidden reference at the same interface.

\begin{samepage}
\begin{center}
\begin{minipage}{\linewidth}
\captionsetup{type=table,hypcap=false}
\caption{\textbf{An example-shared 128-dimensional subspace fails bidirectional effect matching at the initial interfaces.} Gemma/Qwen use visual positions and LLaVA uses answer-adjacent positions, each at its initially selected layer. Values are percentages of the mean clean-minus-masked sequence-score change on the separate 40-example replication set. Energy is the squared hidden-difference norm retained by the projection. Values above 100 indicate effect overshoot, and negative values indicate effect reversal. Rank 128 is the largest tested subspace; no rank from 1 to 128 passes both restoration and corrupted-state-replacement selection gates in any model. Full columns report the unprojected hidden-state reference at the same layer and positions.}
\label{tab:shared-subspace-summary}
\centering
\normalsize
\setlength{\tabcolsep}{2.8pt}
\begin{tabular*}{\linewidth}{@{\extracolsep{\fill}}lrrrrr@{}}
\toprule
Model & Rank-128 energy & Rank-128 restore & Rank-128 delete & Full restore & Full delete \\
\midrule
Gemma & 88.9 & 425.0 & -357.7 & 192.3 & 83.2 \\
Qwen & 53.1 & 53.0 & 51.2 & 94.6 & 76.5 \\
LLaVA & 15.9 & 10.5 & 0.1 & 94.9 & 92.4 \\
\bottomrule
\end{tabular*}
\end{minipage}
\end{center}
\end{samepage}


Retaining substantial hidden-difference energy is not sufficient for causal fidelity. Gemma's rank-128 projection retains 88.9\% of the squared difference but over-restores the answer effect and reverses the deletion effect. Qwen retains 53.1\% of the energy and approximately half of both causal effects. LLaVA retains 15.9\% of the energy, restores 10.5\% of the answer effect, and removes 0.1\%. Failure at every selection rank therefore rejects an example-shared 1--128-dimensional complete subspace at these frozen interfaces. The tested rank range does not rule out higher-rank, nonlinear, or sample-specific representations.

\subsubsection{A result-blind evaluation confirms the LLaVA full-hidden interface}

At the initial frozen interfaces, the unprojected reference gives different results. Gemma's layer-1 visual-position restoration is unstable across the selection and replication sets, while Qwen's layer-14 visual-position replacement effect remains below the 80\% gate. The evaluated LLaVA layer-30 answer-position interface keeps both mean effects within 80--120\% on both sets: restoration and replacement are 93.0\% and 91.0\% on selection and 94.9\% and 92.4\% on replication. We therefore freeze only this LLaVA interface for a separate result-blind test; no layer, position, rank, or threshold is rescanned. Because the three initial interfaces differ in both layer and position type, their cross-model contrast alone cannot assign the difference to model architecture.

The confirmation pool contains 61 masks reviewed without model predictions or intervention results. Human review accepts 60 masks and rejects one before model evaluation. On the 60 accepted examples, the mean clean-minus-masked sequence-score change is $.9378$. Restoring the complete hidden-state difference yields $.8981$, or 95.8\% of that mean change; replacing the clean hidden states with their masked values removes $.8907$, or 95.0\% of the same reference. Percentages use the unrounded means. Table~\ref{tab:llava-hidden-confirmation} shows that all four predeclared paired-bootstrap gates have positive lower bounds.

\begin{samepage}
\begin{center}
\begin{minipage}{\linewidth}
\captionsetup{type=table,hypcap=false}
\caption{\textbf{The frozen LLaVA full-hidden interface passes separate result-blind confirmation.} The interface is layer 30 at the four text positions immediately preceding the answer. Restore and Delete are percentages of the mean clean-minus-masked sequence-score change; Delete writes masked hidden states into the clean run rather than zeroing them. R-L and R-U are raw sequence-log-probability lower bounds for restoration minus 80\% of that mean change and 120\% of that mean change minus restoration; D-L and D-U are the corresponding replacement bounds. All four bounds must be nonnegative, and Gate 1 denotes a pass.}
\label{tab:llava-hidden-confirmation}
\centering
\normalsize
\setlength{\tabcolsep}{2.2pt}
\begin{tabular*}{\linewidth}{@{\extracolsep{\fill}}lrrrrrrrr@{}}
\toprule
Interface & $n$ & Restore & Delete & R-L & R-U & D-L & D-U & Gate \\
\midrule
LLaVA L30, AA-4 & 60 & 95.8 & 95.0 & .050 & .097 & .047 & .100 & 1 \\
\bottomrule
\end{tabular*}
\end{minipage}
\end{center}
\end{samepage}


Restoring the evaluated LLaVA full-hidden interface recovers most of the clean-minus-masked sequence-score difference; replacing its clean value with the masked value removes most of that difference. The 60-example result therefore confirms this interface under both specified interventions. It neither establishes a minimum dimension nor identifies an analogous interface in Gemma or Qwen. The initial comparison's layer-and-position asymmetry motivates one further test before we interpret the cross-model pattern.

\subsubsection{A same-sample late-answer comparison narrows the interface asymmetry}

For this follow-up, each model is tested at its penultimate language layer and the four text positions before the answer, using the same 60 reviewed examples and the same complete-answer score. No result is used to rescan the rule within a model. Table~\ref{tab:matched-late-hidden} reports restoration and corrupted-state replacement on the same image--question IDs. The relative-layer rule was motivated by the already observed LLaVA result, so this comparison is diagnostic; the original LLaVA result above remains its independent confirmation.

\begin{samepage}
\begin{center}
\begin{minipage}{\linewidth}
\captionsetup{type=table,hypcap=false}
\caption{\textbf{Only the evaluated LLaVA late-answer interface passes the matched bidirectional test.} All three models use the same 60 reviewed image--question examples, the zero-based penultimate language layer, four answer-adjacent positions, the same complete-answer score, and the same restored/replaced-state protocol in one runtime. $n_+$ counts examples with a positive clean-minus-masked sequence-score change. $V$ is the mean of that change; $R/V$ and $C/V$ are ratios of within-model means in percent for restoration and corrupted-state replacement. R-L and C-L are raw sequence-score 95\% bootstrap lower bounds for the respective effects minus $0.8V$; the upper bounds must also be nonnegative for Gate=1. Negative ratios denote reversed mean effects. Different models' raw score magnitudes are not compared.}
\label{tab:matched-late-hidden}
\centering
\small
\setlength{\tabcolsep}{2.2pt}
\begin{tabular*}{\linewidth}{@{\extracolsep{\fill}}lrrrrrrrr@{}}
\toprule
Model & Layer & $n_+$ & $V$ & $R/V$ & $C/V$ & R-L & C-L & Gate \\
\midrule
Gemma & 32 & 29/60 & 0.407 & -85.0 & -152.8 & -1.837 & -2.196 & 0 \\
Qwen & 26 & 36/60 & 1.463 & 74.3 & 45.3 & -0.483 & -1.017 & 0 \\
LLaVA & 30 & 41/60 & 0.930 & 95.2 & 95.7 & 0.046 & 0.050 & 1 \\
\bottomrule
\end{tabular*}
\end{minipage}
\end{center}
\end{samepage}


Only the evaluated LLaVA interface passes all four paired bounds in this matched run. Gemma's two mean changes reverse direction; Qwen restores 74.3\% but removes only 45.3\% of its mean clean-minus-masked sequence-score change under corrupted-state replacement. The corresponding LLaVA ratios are 95.2\% and 95.7\%. All 60 examples are scored in every model, but the clean-minus-masked sequence-score change is positive in 29, 36, and 41 examples for Gemma, Qwen, and LLaVA, respectively. Thus, the matched result concerns these frozen interfaces on this LLaVA-origin sample set, not the absence of another successful interface in either Gemma or Qwen. The original LLaVA confirmation and this unified-runtime rerun give slightly different point estimates but the same four-gate decision.

The same late-answer rule is also applied to the 40/20/40 shared-subspace splits. At the matched interfaces, none of the eight ranks from 1 to 128 passes both selection means in any model (Appendix~\ref{app:shared-subspace-details}). This removes the initial visual-versus-answer position difference as an explanation for the tested low-rank failure, while leaving higher-rank and sample-specific objects open. Across the evaluated interfaces, all three models show additional damage from masking evidence rather than matched control content; the tested compact candidates fail their own complete-effect criteria, while one full LLaVA hidden-state interface passes independent bidirectional confirmation against the clean-minus-masked sequence-score change. The Discussion considers what can and cannot be inferred from that separation.

\FloatBarrier
